\chapter{Gaussian vectors and linear Gaussian processes}
\label{chap:pre_gaussian}

This chapter collects the facts about Gaussian vectors in $\R^n$ and about \emph{linear
Gaussian processes} $x \mapsto \ip{x}{G}$ ($G$ a centered Gaussian vector, $x$ ranging over
a compact index set $\Xs \subset \R^d$) that are used throughout Part~\ref{part:main}: closure
of Gaussian vectors under affine maps, identification of the law by mean and covariance,
rotation invariance of pairs of independent standard vectors, independence of uncorrelated
jointly Gaussian vectors, moments of Gaussian norms, tails of empirical means, measurability
and integrability of suprema over compact index sets, the expectation of a maximum of $m$
Gaussians, and comparison by covariance domination (the special case of the
Sudakov--Fernique inequality that suffices for the monotonicity of the Gaussian width,
Lemma~\ref{lem:width_mono}). Everything here is classical
\cite{vershynin2018high, boucheron2013concentration}; the point is to fix the statements at
the level of generality of a library and to name what exists.

\textbf{What Mathlib provides.} The standard Gaussian measure on a finite-dimensional inner
product space, \texttt{ProbabilityTheory.stdGaussian}, and the multivariate Gaussian
\texttt{multivariateGaussian μ S} on \texttt{EuclideanSpace ℝ ι}, defined as the image of
\texttt{stdGaussian} under $x \mapsto \mu + S^{1/2}x$ with $S^{1/2} = $ \texttt{CFC.sqrt S}
(\texttt{Mathlib/Probability/Distributions/Gaussian/Multivariate.lean}), with
\texttt{integral\_id\_multivariateGaussian}, \texttt{covarianceBilin\_multivariateGaussian},
\texttt{charFun\_multivariateGaussian}, \texttt{multivariateGaussian\_zero\_one},
\texttt{stdGaussian\_map} (invariance under linear isometry equivalences),
\texttt{map\_pi\_eq\_stdGaussian} and \texttt{stdGaussian\_eq\_map\_pi\_orthonormalBasis}
(coordinates are i.i.d.\ standard). The class \texttt{IsGaussian} of Gaussian measures
(\texttt{Gaussian/Basic.lean}: every continuous linear functional has a real Gaussian law),
closed under continuous linear maps (\texttt{isGaussian\_map}), translations, negation and
products, characterized by its characteristic function
(\texttt{IsGaussian.charFun\_eq}, \texttt{isGaussian\_iff\_charFun\_eq}); the predicate
\texttt{HasGaussianLaw X P} on random variables
(\texttt{Gaussian/HasGaussianLaw/Basic.lean}: \texttt{HasGaussianLaw.map}, \texttt{.smul},
\texttt{.neg}, \texttt{.add}, \texttt{.eval}, \texttt{.integrable}, \texttt{.memLp}), and the
independence results of \texttt{Gaussian/HasGaussianLaw/Independence.lean}
(\texttt{iIndepFun.hasGaussianLaw}, \texttt{IndepFun.hasGaussianLaw},
\texttt{HasGaussianLaw.indepFun\_of\_covariance\_eval}, \texttt{\_inner},
\texttt{\_eq\_zero}). Fernique's theorem \texttt{IsGaussian.exists\_integrable\_exp\_sq}
(\texttt{Gaussian/Fernique.lean}) and the rotation invariance
\texttt{IsGaussian.map\_rotation\_eq\_self} of $\mu \otimes \mu$ for centered $\mu$. On the real
line, \texttt{gaussianReal}, \texttt{mgf\_gaussianReal}, \texttt{integral\_id\_gaussianReal},
\texttt{variance\_id\_gaussianReal}, \texttt{gaussianReal\_map\_const\_mul},
\texttt{gaussianReal\_add\_gaussianReal\_of\_indepFun} (\texttt{Gaussian/Real.lean}).
Sub-Gaussian tails: \texttt{HasSubgaussianMGF} with \texttt{measure\_ge\_le}
(\texttt{Mathlib/Probability/Moments/SubGaussian.lean}) and \texttt{measure\_le\_le}
(LML, \texttt{ForMathlib/Probability/Moments/SubGaussian.lean}).

\textbf{What is missing.} The matrix form of the linear-image lemma (mean and covariance
matrix of $b + LV$), the uniqueness of the law given mean and covariance matrix in the
matrix language of Part~\ref{part:main}, the fourth moment of a Gaussian norm, the
lower bound $\E\norm{Y} \ge \sqrt{\tr S}/\sqrt 3$, the expectation bound for a maximum of $m$
Gaussians, the countable-dense reduction for suprema over compact index sets, and the
covariance-domination comparison lemma. All are short.

\textbf{Lean remark.} Vectors live in \texttt{EuclideanSpace ℝ (Fin d)} and matrices in
\texttt{Matrix (Fin d) (Fin d) ℝ}, acting through \texttt{Matrix.toEuclideanLin} /
\texttt{toEuclideanCLM}. The positive semidefinite square root of $S \in \PSD$ is
\texttt{CFC.sqrt S}, the one used by \texttt{multivariateGaussian}; for $S \in \PD$ it agrees
with the square root of Lemma~\ref{lem:sqrt_props}. The random vectors of this chapter are
functions \texttt{Ω → EuclideanSpace ℝ (Fin d)} with \texttt{HasLaw} a given Gaussian measure,
or with \texttt{HasGaussianLaw} together with prescribed mean and covariance.

\section{Gaussian vectors}
\label{sec:pre_gaussian_vectors}

\begin{definition}[Gaussian measures and Gaussian vectors]\label{def:pg_gaussian_vector}
Let $n \ge 1$, $\mu \in \R^n$ and $S \in \PSD$ (an $n \times n$ positive semidefinite
matrix). The \emph{Gaussian measure} $\Normal(\mu, S)$ on $\R^n$ is the image of the standard
Gaussian measure $\Normal(0, I_n)$ (Mathlib \texttt{stdGaussian}) under
$x \mapsto \mu + S^{1/2}x$, where $S^{1/2}$ is the positive semidefinite square root of $S$
(Mathlib \texttt{multivariateGaussian μ S}). A random vector $V : \Omega \to \R^n$ satisfies
$V \sim \Normal(\mu, S)$ if its law is $\Normal(\mu,S)$. More generally $V$ is a
\emph{Gaussian vector} if $\ip{t}{V}$ has a (possibly degenerate) real Gaussian law for every
$t \in \R^n$ (Mathlib \texttt{HasGaussianLaw V P}); its \emph{mean} is $\E V$ and its
\emph{covariance matrix} is $\operatorname{Cov}(V) := \E[(V - \E V)(V - \E V)^\top]$.
For $n = 1$, $\Normal(\mu, \sigma^2)$ is Mathlib's \texttt{gaussianReal μ σ²}; for
$\sigma = 0$ it is the Dirac mass at $\mu$.
\end{definition}

We recall the facts that Mathlib attaches to this definition: $\Normal(\mu,S)$ is a
probability measure with mean $\mu$ (\texttt{integral\_id\_multivariateGaussian}), covariance
bilinear form $(s,t) \mapsto s^\top S t$ (\texttt{covarianceBilin\_multivariateGaussian}, for
$S \in \PSD$), and characteristic function
$t \mapsto \exp(i\ip{t}{\mu} - \tfrac12 t^\top S t)$ (\texttt{charFun\_multivariateGaussian});
$\Normal(0, I_n) = $ \texttt{stdGaussian} (\texttt{multivariateGaussian\_zero\_one}); every
Gaussian vector is integrable with all moments (\texttt{HasGaussianLaw.memLp}).

\begin{lemma}[Standard Gaussian vectors have i.i.d.\ coordinates]\label{lem:gauss_pi}
\lean{ProbabilityTheory.hasLaw_pi_of_hasCondDistrib_const, ProbabilityTheory.iIndepFun_of_hasCondDistrib_const, ProbabilityTheory.hasLaw_of_hasCondDistrib_const}\leanok
\uses{def:pg_gaussian_vector}
Let $T \ge 1$.
\begin{enumerate}
  \item[(i)] If $\eta_0, \dots, \eta_{T-1}$ are independent real random variables, each with
        law $\Normal(0,1)$, then the random vector $\eta := (\eta_t)_{t<T} \in \R^T$
        satisfies $\eta \sim \Normal(0, I_T)$.
  \item[(ii)] Conversely, if $\eta \sim \Normal(0, I_T)$ then its coordinates
        $\eta_0,\dots,\eta_{T-1}$ are independent with law $\Normal(0,1)$.
\end{enumerate}
\end{lemma}

\begin{proof}
\uses{def:pg_gaussian_vector}
(i) The joint law of independent random variables is the product of their laws
(Mathlib \texttt{iIndepFun\_iff\_map\_fun\_eq\_pi\_map}), so the law of $\eta$, seen in
$\R^T = $ \texttt{EuclideanSpace ℝ (Fin T)}, is the image of
$\bigotimes_{t<T}\Normal(0,1)$ under the identification \texttt{WithLp.toLp 2}, which is
\texttt{stdGaussian} by \texttt{map\_pi\_eq\_stdGaussian}; conclude with
\texttt{multivariateGaussian\_zero\_one}. (ii) The same identity read backwards through the
measurable equivalence \texttt{WithLp.toLp 2}: the law of $(\eta_t)_t$ as an element of
$\mathrm{Fin}\,T \to \R$ is the product measure, which is the independence of the
coordinates with the stated marginals (\texttt{iIndepFun\_iff\_map\_fun\_eq\_pi\_map} again,
or \texttt{measurePreserving\_eval\_multivariateGaussian} for the marginals).
\end{proof}

\textbf{Lean remark.} Part (i) is Mathlib's \texttt{iIndepFun.hasLaw\_pi} together with
\texttt{map\_pi\_eq\_stdGaussian}, and (ii) is \texttt{iIndepFun\_iff\_hasLaw\_pi\_pi}. The form
in which (i) is used is sequential: if $\eta_0 \sim \nu$ and, for every $n$, the conditional law
of $\eta_{n+1}$ given $(\eta_0,\dots,\eta_n)$ (or given any variable of which this vector is a
measurable function) is the constant $\nu$, then $(\eta_0,\dots,\eta_{n-1})$ has the product law
$\nu^{\otimes n}$ and the $\eta_t$ are independent with law $\nu$
(\texttt{ProbabilityTheory.hasLaw\_pi\_of\_hasCondDistrib\_const},
\texttt{iIndepFun\_of\_hasCondDistrib\_const}; proof by induction through the measurable
equivalence \texttt{MeasurableEquiv.piFinSuccAbove}). This is what
Lemma~\ref{lem:noise_representation} provides for the noise of a linear Gaussian run.

\begin{lemma}[The law of a Gaussian vector is determined by its mean and covariance]
\label{lem:gauss_law_of_cov}
\uses{def:pg_gaussian_vector}
Let $V, V'$ be Gaussian vectors in $\R^n$ (possibly on different probability spaces).
\begin{enumerate}
  \item[(i)] If $\E V = \E V'$ and $\operatorname{Cov}(V) = \operatorname{Cov}(V')$, then $V$
        and $V'$ have the same law.
  \item[(ii)] If $\E V = \mu$ and $\operatorname{Cov}(V) = S$, then $V \sim \Normal(\mu, S)$.
  \item[(iii)] For $S \in \PSD$ and $g \sim \Normal(0, I_n)$, $S^{1/2} g \sim \Normal(0,S)$.
  \item[(iv)] If $G \sim \Normal(0, S)$ then $-G \sim \Normal(0,S)$, i.e.\ $-G \eqd G$.
\end{enumerate}
\end{lemma}

\begin{proof}
\uses{def:pg_gaussian_vector}
(i) The characteristic function of a Gaussian vector is
$t \mapsto \exp\big(i\,\E\ip{t}{V} - \tfrac12\Var\ip{t}{V}\big)$
(Mathlib \texttt{IsGaussian.charFun\_eq}, \texttt{HasGaussianLaw.charFun\_map\_eq}), and
$\E\ip{t}{V} = \ip{t}{\E V}$, $\Var\ip{t}{V} = t^\top\operatorname{Cov}(V)\,t$ (bilinearity of
the covariance, Mathlib \texttt{covarianceBilin}). Hence $V$ and $V'$ have the same
characteristic function, and a finite measure on a finite-dimensional inner product space is
determined by its characteristic function (\texttt{Measure.ext\_of\_charFun}). (ii) The
measure $\Normal(\mu,S)$ is Gaussian (\texttt{isGaussian\_multivariateGaussian}) with mean
$\mu$ and covariance $S$ (the facts recalled after Definition~\ref{def:pg_gaussian_vector});
apply (i) to $V$ and to the identity on $(\R^n, \Normal(\mu,S))$. (iii) is the definition of
$\Normal(0,S)$ (Mathlib \texttt{multivariateGaussian}, as the image of \texttt{stdGaussian}),
using that the law of $S^{1/2}g$ is the image of the law of $g$. (iv) $-G$ is a Gaussian
vector (\texttt{HasGaussianLaw.neg}) with mean $0$ and covariance
$\E[(-G)(-G)^\top] = S$; apply (ii).
\end{proof}

\begin{lemma}[Affine images of Gaussian vectors]\label{lem:gauss_linear_image}
\lean{ProbabilityTheory.multivariateGaussian_map_toEuclideanCLM, ProbabilityTheory.multivariateGaussian_map_toEuclideanLin}
\uses{def:pg_gaussian_vector}
Let $V \sim \Normal(\mu, S)$ in $\R^n$, $L \in \R^{m \times n}$ and $b \in \R^m$. Then
$b + LV \sim \Normal(b + L\mu,\ L S L^\top)$. In particular, for $g \sim \Normal(0, I_n)$,
$M \in \R^{m\times n}$ and $\mu \in \R^m$: $\mu + Mg \sim \Normal(\mu, MM^\top)$; and for
$a \in \R^n$, $\ip{a}{g} \sim \Normal(0, \norm{a}^2)$.
\end{lemma}

\begin{proof}
\uses{def:pg_gaussian_vector, lem:gauss_law_of_cov}
$b + LV$ is a Gaussian vector: $V$ is one, $L$ is a continuous linear map
(\texttt{HasGaussianLaw.map}) and translations preserve Gaussianity
(\texttt{IsGaussian} instance for \texttt{μ.map (· + c)}). Its mean is $b + L\mu$ by
linearity of the integral, and for $s, t \in \R^m$,
$\operatorname{Cov}(\ip{s}{b + LV}, \ip{t}{b+LV}) = \operatorname{Cov}(\ip{L^\top s}{V},
\ip{L^\top t}{V}) = (L^\top s)^\top S (L^\top t) = s^\top (LSL^\top) t$, i.e.\ the covariance
matrix is $LSL^\top$. Conclude with Lemma~\ref{lem:gauss_law_of_cov}(ii). The special cases
are $S = I_n$ and $L = a^\top$ (with $m = 1$, $\Normal(0, a^\top a)$ on $\R$ being
\texttt{gaussianReal 0 ‖a‖²}).
\end{proof}

\textbf{Lean remark.} With $V$ given by \texttt{HasLaw V (multivariateGaussian μ S) P}, the
conclusion is \texttt{HasLaw (fun ω ↦ b + L (V ω)) (multivariateGaussian (b + L μ) (L * S * Lᵀ)) P}.
The proof through characteristic functions is an alternative to the one through
\texttt{HasGaussianLaw}: $\E e^{i\ip{t}{b + LV}} = e^{i\ip{t}{b}}\,\E e^{i\ip{L^\top t}{V}}$
and \texttt{charFun\_multivariateGaussian} on both sides.

\begin{lemma}[Independent Gaussian vectors are jointly Gaussian]\label{lem:pg_indep_gaussian_pair}
\lean{ProbabilityTheory.multivariateGaussian_conv}
\uses{def:pg_gaussian_vector}
Let $V_1 \sim \Normal(\mu_1, S_1)$ in $\R^{n_1}$ and $V_2 \sim \Normal(\mu_2, S_2)$ in
$\R^{n_2}$ be independent, on the same probability space. Then the pair $(V_1, V_2)$, seen in
$\R^{n_1 + n_2}$, satisfies
$(V_1,V_2) \sim \Normal\big((\mu_1,\mu_2), \operatorname{diag}(S_1, S_2)\big)$.
If $n_1 = n_2$ then $V_1 + V_2 \sim \Normal(\mu_1 + \mu_2, S_1 + S_2)$.
\end{lemma}

\begin{proof}
\uses{def:pg_gaussian_vector, lem:gauss_law_of_cov, lem:gauss_linear_image}
The pair is a Gaussian vector (Mathlib \texttt{IndepFun.hasGaussianLaw}: the law of the pair
is the product measure, and a product of Gaussian measures is Gaussian). Its mean is
$(\mu_1,\mu_2)$ and its covariance matrix is block diagonal: the cross covariances
$\operatorname{Cov}(\ip{s}{V_1}, \ip{t}{V_2})$ vanish because $\ip{s}{V_1}$ and $\ip{t}{V_2}$
are independent and square-integrable (the covariance of independent variables vanishes:
Mathlib \texttt{IndepFun.covariance\_eq\_zero} in \texttt{Mathlib/Probability/Moments/Covariance.lean},
or the product formula \texttt{IndepFun.integral\_mul\_eq\_mul\_integral}). Conclude with
Lemma~\ref{lem:gauss_law_of_cov}(ii). For the sum apply Lemma~\ref{lem:gauss_linear_image}
with $L = [I_n\ I_n]$: $L\operatorname{diag}(S_1,S_2)L^\top = S_1 + S_2$.
\end{proof}

\begin{lemma}[Orthogonal invariance and rotation of pairs]\label{lem:gauss_rotation}
\uses{def:pg_gaussian_vector}
\begin{enumerate}
  \item[(i)] If $g \sim \Normal(0, I_d)$ and $U \in \R^{d\times d}$ is orthogonal
        ($U^\top U = I_d$), then $Ug \sim \Normal(0, I_d)$.
  \item[(ii)] If $g, g'$ are independent, each with law $\Normal(0, I_d)$, then for every
        $\theta \in \R$ the pair
        \[
          (g_\theta, g'_\theta) := \big(g\sin\theta + g'\cos\theta,\ g\cos\theta - g'\sin\theta\big)
        \]
        has the same law as $(g, g')$, namely $\Normal(0, I_{2d})$ on $\R^d \times \R^d$;
        in particular $g_\theta$ and $g'_\theta$ are independent $\Normal(0,I_d)$ vectors.
\end{enumerate}
\end{lemma}

\begin{proof}
\uses{def:pg_gaussian_vector, lem:pg_indep_gaussian_pair, lem:gauss_pi, lem:gauss_uncorrelated_indep}
(i) $U$ is a linear isometry equivalence of $\R^d$ and the standard Gaussian measure is
invariant under linear isometry equivalences (Mathlib \texttt{stdGaussian\_map}); the law of
$Ug$ is the image of the law of $g$. (ii) By Lemma~\ref{lem:pg_indep_gaussian_pair} the pair
$(g,g')$ has law $\Normal(0, I_{2d})$. The map
$R_\theta(u,v) := (u\sin\theta + v\cos\theta,\ u\cos\theta - v\sin\theta)$ is linear and
\[
  \norm{u\sin\theta + v\cos\theta}^2 + \norm{u\cos\theta - v\sin\theta}^2
  = \norm{u}^2(\sin^2\theta + \cos^2\theta) + \norm{v}^2(\cos^2\theta + \sin^2\theta)
    + 2\ip{u}{v}(\sin\theta\cos\theta - \cos\theta\sin\theta) = \norm{u}^2 + \norm{v}^2,
\]
so $R_\theta$ is a linear isometry of $\R^{2d}$ (bijective since it is an isometry of a
finite-dimensional space, or explicitly $R_\theta^{-1} = R_\theta^\top$). Apply (i) in
dimension $2d$. (Mathlib's \texttt{IsGaussian.map\_rotation\_eq\_self} is the same statement
for the map \texttt{ContinuousLinearMap.rotation θ} of
\texttt{Mathlib/Probability/Distributions/Fernique.lean}, namely
$(u,v)\mapsto(u\cos\theta + v\sin\theta,\ -u\sin\theta + v\cos\theta)$, applied to
$\mu\otimes\mu$ with $\mu$ a centered Gaussian measure. Since
$\mathrm{rotation}(-\theta)(u,v) = (u\cos\theta - v\sin\theta,\ u\sin\theta + v\cos\theta)$,
one has $R_\theta = \mathrm{swap}\circ\mathrm{rotation}(-\theta)$ with
$\mathrm{swap}(u,v) := (v,u)$, and the swap also preserves $\mu\otimes\mu$
(\texttt{Measure.prod\_swap}).) Independence of the components of
a $\Normal(0,I_{2d})$ vector is Lemma~\ref{lem:gauss_pi}(ii) applied blockwise, or
Lemma~\ref{lem:gauss_uncorrelated_indep}.
\end{proof}

\begin{lemma}[Joint law of two linear images]\label{lem:gauss_joint_linear}
\uses{def:pg_gaussian_vector}
Let $V \sim \Normal(\mu,S)$ in $\R^n$ and $L_1 \in \R^{m_1\times n}$, $L_2 \in \R^{m_2\times n}$.
Then $(L_1 V, L_2 V)$ is a Gaussian vector in $\R^{m_1}\times\R^{m_2}$ with
\[
  (L_1V, L_2V) \sim \Normal\Big( (L_1\mu, L_2\mu),
  \begin{pmatrix} L_1SL_1^\top & L_1SL_2^\top \\ L_2SL_1^\top & L_2SL_2^\top\end{pmatrix}\Big),
\]
and in particular the cross covariance is
$\operatorname{Cov}(L_1V, L_2V) := \E[(L_1V - L_1\mu)(L_2V - L_2\mu)^\top] = L_1 S L_2^\top$.
\end{lemma}

\begin{proof}
\uses{def:pg_gaussian_vector, lem:gauss_linear_image}
$(L_1V, L_2V) = LV$ with $L := \binom{L_1}{L_2} \in \R^{(m_1+m_2)\times n}$;
Lemma~\ref{lem:gauss_linear_image} gives $LV \sim \Normal(L\mu, LSL^\top)$ and
$LSL^\top$ is the displayed block matrix. The cross covariance is its off-diagonal block.
\end{proof}

\begin{lemma}[Uncorrelated jointly Gaussian vectors are independent]\label{lem:gauss_uncorrelated_indep}
\uses{def:pg_gaussian_vector}
Let $(V_1, V_2)$ be a Gaussian vector in $\R^{m_1}\times\R^{m_2}$ (i.e.\ the pair is jointly
Gaussian) such that $\operatorname{Cov}(V_1, V_2) = 0$, i.e.\
$\operatorname{Cov}((V_1)_i, (V_2)_j) = 0$ for all $i \le m_1$, $j \le m_2$. Then $V_1$ and
$V_2$ are independent.
\end{lemma}

\begin{proof}
\uses{def:pg_gaussian_vector}
Mathlib: \texttt{HasGaussianLaw.indepFun\_of\_covariance\_eval} (coordinates) or
\texttt{HasGaussianLaw.indepFun\_of\_covariance\_inner} (all $\operatorname{Cov}(\ip{s}{V_1},
\ip{t}{V_2}) = 0$, which follows from the coordinatewise hypothesis by bilinearity).
\end{proof}

\begin{lemma}[Expectation of an inner product with an independent centered vector]
\label{lem:indep_integral_zero}
Let $W : \Omega \to \R^d$ be integrable with $\E W = 0$, let $V : \Omega \to \mathcal V$ be a
random variable with values in a measurable space, independent of $W$, and let
$f : \mathcal V \to \R^d$ be bounded and measurable. Then $\ip{f(V)}{W}$ is integrable and
$\E\ip{f(V)}{W} = 0$.
\end{lemma}

\begin{proof}
$\ip{f(V)}{W} = \sum_{i<d} f_i(V)\,W_i$. For each $i$, $f_i(V)$ is bounded and measurable and
$W_i$ is integrable, so their product is integrable; $f_i(V)$ and $W_i$ are independent
(compositions of independent variables with measurable maps, Mathlib \texttt{IndepFun.comp}),
so $\E[f_i(V)W_i] = \E[f_i(V)]\,\E[W_i] = 0$ (product formula for independent integrable
variables, Mathlib \texttt{IndepFun.integral\_mul\_eq\_mul\_integral}). Sum over $i$.
\end{proof}

\section{Moments and tails}
\label{sec:pre_gaussian_moments}

\begin{lemma}[Moments of a centered real Gaussian]\label{lem:gauss_1d_moments}
\uses{def:pg_gaussian_vector}
Let $\sigma \ge 0$ and $Z \sim \Normal(0,\sigma^2)$. Then $Z$ has finite moments of all
orders and
\[
  \E Z = 0,\quad \E Z^2 = \sigma^2,\quad \E Z^4 = 3\sigma^4,\quad
  \E\abs{Z} = \sigma\sqrt{2/\pi},\quad \E e^{\lambda Z} = e^{\lambda^2\sigma^2/2}\ (\lambda\in\R).
\]
\end{lemma}

\begin{proof}
\uses{def:pg_gaussian_vector}
Mathlib: \texttt{memLp\_id\_gaussianReal} (all moments), \texttt{integral\_id\_gaussianReal},
\texttt{variance\_id\_gaussianReal}, \texttt{mgf\_gaussianReal} (with mean $0$ and variance
$\sigma^2$: $\E e^{\lambda Z} = \exp(0\cdot\lambda + \sigma^2\lambda^2/2)$). For the fourth
moment and the absolute moment, reduce to $\sigma = 1$ by $Z \eqd \sigma\xi$,
$\xi \sim \Normal(0,1)$ (\texttt{gaussianReal\_map\_const\_mul}; for $\sigma = 0$ both sides
vanish), and compute with the density $\varphi(x) = e^{-x^2/2}/\sqrt{2\pi}$:
$\int x^4\varphi = 3\int x^2\varphi = 3$ by integration by parts
($x^3\cdot x\varphi(x) = -x^3\varphi'(x)$, \texttt{integral\_mul\_deriv\_eq\_deriv\_mul}), and
$\int\abs{x}\varphi = 2\int_0^\infty x\varphi(x)\,dx = 2\varphi(0) = \sqrt{2/\pi}$ since
$x\varphi(x) = -\varphi'(x)$. Alternatively, $\E Z^4$ is the fourth derivative at $0$ of the
MGF $e^{\lambda^2\sigma^2/2}$ (Mathlib \texttt{iteratedDeriv\_mgf\_zero}).
\end{proof}

\begin{lemma}[Moments of the norm of a Gaussian vector]\label{lem:gauss_norm_moments}
\lean{ProbabilityTheory.integral_norm_sq_multivariateGaussian, ProbabilityTheory.integral_norm_le_sqrt_trace_multivariateGaussian, ProbabilityTheory.integral_norm_stdGaussian_le}
\uses{def:pg_gaussian_vector}
Let $S \in \PSD$ ($d \times d$) and $Y \sim \Normal(0, S)$; equivalently $Y = S^{1/2}g$ with
$g \sim \Normal(0,I_d)$. Then
\begin{enumerate}
  \item[(i)] $\E\norm{Y}^2 = \tr S$;
  \item[(ii)] $\E\norm{Y}^4 = (\tr S)^2 + 2\tr(S^2) \le 3(\tr S)^2$;
  \item[(iii)] $\E\norm{Y} \le \sqrt{\tr S}$;
  \item[(iv)] $\E\norm{Y} \ge (\tr S)^{3/2}\big/\sqrt{(\tr S)^2 + 2\tr(S^2)} \ge \sqrt{\tr S}/\sqrt3$.
\end{enumerate}
In the special case $S = I_d$: $\E\norm{g}^2 = d$, $\E\norm{g}^4 = d^2 + 2d$,
$\E\norm{g} \le \sqrt d$ and $\E\norm{g} \ge d/\sqrt{d+2} \ge \sqrt{d/3}$.
\end{lemma}

\begin{proof}
\uses{def:pg_gaussian_vector, lem:gauss_law_of_cov, lem:gauss_rotation, lem:gauss_pi, lem:gauss_1d_moments}
All quantities depend only on the law of $Y$, so by Lemma~\ref{lem:gauss_law_of_cov}(iii) we
may take $Y = S^{1/2}g$. Diagonalize $S = U\operatorname{diag}(s_1,\dots,s_d)U^\top$ with $U$
orthogonal and $s_i \ge 0$ (spectral theorem for symmetric matrices, Mathlib
\texttt{Matrix.IsHermitian.spectral\_theorem}; nonnegativity of the eigenvalues from
$S \in \PSD$). Then $\norm{Y}^2 = g^\top S g = \sum_i s_i h_i^2$ with $h := U^\top g$, and
$h \sim \Normal(0,I_d)$ by Lemma~\ref{lem:gauss_rotation}(i), so the $h_i$ are i.i.d.\
$\Normal(0,1)$ (Lemma~\ref{lem:gauss_pi}(ii)). (i) $\E\sum_i s_ih_i^2 = \sum_i s_i = \tr S$
(Lemma~\ref{lem:gauss_1d_moments}). (ii) Expanding the square and using independence,
$\E h_i^4 = 3$ and $\E h_i^2 = 1$,
\[
  \E\Big(\sum_i s_ih_i^2\Big)^2 = \sum_i s_i^2\,\E h_i^4 + \sum_{i\ne j}s_is_j\,\E h_i^2\,\E h_j^2
  = 3\sum_i s_i^2 + \Big(\big(\sum_i s_i\big)^2 - \sum_i s_i^2\Big) = (\tr S)^2 + 2\tr(S^2),
\]
using $\tr(S^2) = \sum_i s_i^2$. Since $s_i \ge 0$, $\sum_i s_i^2 \le (\sum_i s_i)^2$, whence
the bound $3(\tr S)^2$. (iii) is Jensen (or Cauchy--Schwarz): $(\E\norm{Y})^2 \le \E\norm{Y}^2$.
(iv) Hölder's inequality with exponents $3/2$ and $3$ applied to
$\norm{Y}^2 = \norm{Y}^{2/3}\cdot\norm{Y}^{4/3}$ gives
$\E\norm{Y}^2 \le (\E\norm{Y})^{2/3}(\E\norm{Y}^4)^{1/3}$, i.e.\
$(\E\norm{Y})^{2/3} \ge \tr S\,/\,((\tr S)^2 + 2\tr(S^2))^{1/3}$; raise to the power $3/2$ and
use (ii) for the last inequality. For $S = I_d$, $\tr S = d$, $\tr(S^2) = d$ and
$d^{3/2}/\sqrt{d^2+2d} = d/\sqrt{d+2}$, and $d/\sqrt{d+2} \ge \sqrt{d/3}$ since
$d \ge 1$ gives $d \ge (d+2)/3$.
\end{proof}

\begin{lemma}[Gaussian tail]\label{lem:gauss_tail}
\uses{def:pg_gaussian_vector}
Let $\mu \in \R$, $\sigma > 0$ and $Z \sim \Normal(\mu,\sigma^2)$. Then $Z - \mu$ has a
sub-Gaussian moment generating function with constant $\sigma^2$
(Mathlib \texttt{HasSubgaussianMGF (Z - μ) σ²}: $\E e^{\lambda(Z-\mu)} \le e^{\sigma^2\lambda^2/2}$
for all $\lambda$, with equality), and for every $t \ge 0$,
\[
  \Prob(Z - \mu \ge t) \le e^{-t^2/(2\sigma^2)}, \qquad
  \Prob(Z - \mu \le -t) \le e^{-t^2/(2\sigma^2)}, \qquad
  \Prob(\abs{Z - \mu} \ge t) \le 2e^{-t^2/(2\sigma^2)} .
\]
\end{lemma}

\begin{proof}
\uses{def:pg_gaussian_vector}
$Z - \mu \sim \Normal(0,\sigma^2)$ (\texttt{gaussianReal\_sub\_const}) and
\texttt{mgf\_gaussianReal} gives $\E e^{\lambda(Z-\mu)} = e^{\sigma^2\lambda^2/2}$, with
integrability from \texttt{integrable\_exp\_mul\_gaussianReal}; this is the structure
\texttt{HasSubgaussianMGF}. The tails are the Chernoff bounds
\texttt{HasSubgaussianMGF.measure\_ge\_le} (Mathlib) and
\texttt{HasSubgaussianMGF.measure\_le\_le} (LML); the two-sided bound is a union bound.
\end{proof}

\begin{lemma}[Concentration of a Gaussian empirical mean]\label{lem:gauss_mean_concentration}
\uses{def:pg_gaussian_vector, lem:gauss_tail}
Let $n \ge 1$ and $\eta_1,\dots,\eta_n$ be i.i.d.\ $\Normal(0,1)$. Then
$\bar\eta := \frac1n\sum_{i\le n}\eta_i \sim \Normal(0, 1/n)$ and for all $t \ge 0$,
$\Prob(\abs{\bar\eta} \ge t) \le 2e^{-nt^2/2}$. Consequently, for $\eps > 0$ and
$\delta \in (0,1)$, if $n \ge 8\log(2/\delta)/\eps^2$ then
$\Prob(\abs{\bar\eta} \le \eps/2) \ge 1 - \delta$. The same holds for $\bar y - \mu$ when
$y_i = \mu + \eta_i$.
\end{lemma}

\begin{proof}
\uses{lem:gauss_tail, lem:gauss_pi, lem:gauss_linear_image}
$\sum_i\eta_i \sim \Normal(0,n)$ (sums of independent real Gaussians, Mathlib
\texttt{gaussianReal\_add\_gaussianReal\_of\_indepFun} by induction; or
Lemmas~\ref{lem:gauss_pi} and~\ref{lem:gauss_linear_image} with $a = (1,\dots,1)$), and
$\bar\eta \sim \Normal(0, 1/n)$ (\texttt{gaussianReal\_map\_const\_mul}). Lemma~\ref{lem:gauss_tail}
with $\sigma^2 = 1/n$ gives $\Prob(\abs{\bar\eta}\ge t) \le 2e^{-t^2/(2/n)} = 2e^{-nt^2/2}$.
With $t = \eps/2$: $2e^{-n\eps^2/8} \le 2e^{-\log(2/\delta)} = \delta$ when
$n\eps^2/8 \ge \log(2/\delta)$.
\end{proof}

\section{Suprema of linear Gaussian processes}
\label{sec:pre_gaussian_sup}

\begin{lemma}[Suprema over a compact set reduce to a countable dense subset]
\label{lem:sup_countable_dense}
\lean{supportFn, lipschitzWith_supportFn, convexOn_supportFn, continuous_supportFn, abs_supportFn_le, exists_supportFn_eq_inner, partialSupInner, monotone_partialSupInner, abs_partialSupInner_le, exists_forall_supportFn_eq_iSup, tendsto_partialSupInner}\leanok
Let $\Xs \subset \R^d$ be nonempty and compact and $R := \max_{x\in\Xs}\norm{x}$.
\begin{enumerate}
  \item[(i)] There is a countable dense subset $D \subseteq \Xs$; fix an enumeration
        $D = \{q_n : n \in \N\}$ (repetitions allowed).
  \item[(ii)] For every $v \in \R^d$ the supremum $m(v) := \sup_{x\in\Xs}\ip{x}{v}$ is
        attained (it is a maximum), and
        $m(v) = \sup_{n\in\N}\ip{q_n}{v} = \lim_{n\to\infty}\max_{i\le n}\ip{q_i}{v}$,
        the sequence $\max_{i\le n}\ip{q_i}{v}$ being nondecreasing in $n$.
  \item[(iii)] $m : \R^d \to \R$ is convex, $R$-Lipschitz ($\abs{m(v) - m(v')} \le R\norm{v - v'}$),
        hence continuous and Borel measurable, and $\abs{m(v)} \le R\norm{v}$.
  \item[(iv)] For every random vector $G : \Omega \to \R^d$, $\sup_{x\in\Xs}\ip{x}{G} = m(G)$ is
        a real random variable with $\abs{m(G)} \le R\norm{G}$; it is integrable as soon as
        $\norm{G}$ is.
\end{enumerate}
\end{lemma}

\begin{proof}
\leanok
(i) $\Xs$ is a separable metric space (Mathlib \texttt{TopologicalSpace.exists\_countable\_dense}
for the subtype, which is second countable as a subspace of $\R^d$). (ii) For fixed $v$ the
map $x \mapsto \ip{x}{v}$ is continuous, so it attains its maximum on the compact $\Xs$
(\texttt{IsCompact.exists\_isMaxOn}). Clearly $\sup_{n}\ip{q_n}{v} \le m(v)$. Conversely, if
$x \in \Xs$ attains $m(v)$, density gives $q_{n_k} \to x$ and continuity gives
$\ip{q_{n_k}}{v} \to \ip{x}{v} = m(v)$, so $\sup_n\ip{q_n}{v} \ge m(v)$. The partial maxima are
nondecreasing and converge to the supremum of the sequence
(\texttt{tendsto\_atTop\_ciSup} for monotone bounded sequences). (iii) $m$ is a supremum of
linear functions, hence convex. For $v, v'$ and $x \in \Xs$ attaining $m(v)$,
$m(v) - m(v') \le \ip{x}{v} - \ip{x}{v'} \le \norm{x}\norm{v-v'} \le R\norm{v-v'}$; symmetrize.
$m(0) = 0$ gives $\abs{m(v)} \le R\norm{v}$. Measurability follows from continuity, or
directly from (ii): a countable supremum of measurable functions is measurable
(\texttt{Measurable.iSup}). (iv) is (ii)--(iii) composed with $G$.
\end{proof}

\textbf{Lean remark.} In Lean, $\sup_{x\in\Xs}\ip{x}{v}$ is best written as
\texttt{⨆ x : Xs, ⟪(x : E), v⟫} (a \texttt{ciSup} over the compact subtype, bounded above),
and item (ii) provides the equality with \texttt{⨆ n, ⟪q n, v⟫} for a sequence
\texttt{q : ℕ → Xs} with dense range (\texttt{TopologicalSpace.exists\_dense\_seq}). All
measure-theoretic statements about $\sup_{x\in\Xs}\ip{x}{G}$ in this blueprint are proved
through this countable form, and Chapter~\ref{chap:pre_concentration} uses the partial maxima
$\max_{i\le n}\ip{q_i}{G}$ as the finite approximations.

\begin{lemma}[Expected maximum of $m$ sub-Gaussian variables]\label{lem:gauss_max_upper}
\uses{lem:gauss_tail}
Let $m \ge 1$, $\sigma \ge 0$ and $Z_1,\dots,Z_m$ be real random variables on a common
probability space, each with a sub-Gaussian moment generating function with constant
$\sigma^2$ (Mathlib \texttt{HasSubgaussianMGF (Z i) σ²}: $e^{\beta Z_i}$ is integrable and
$\E e^{\beta Z_i}\le e^{\beta^2\sigma^2/2}$ for every $\beta\in\R$; no assumption on their
joint law). Then $\max_{i\le m}Z_i$ is integrable and
\[
  \E\max_{i\le m}Z_i \le \sigma\sqrt{2\log m}.
\]
In particular this applies to centered Gaussian variables $Z_i \sim \Normal(0,\sigma_i^2)$
with $\sigma_i^2 \le \sigma^2$ for every $i$, whatever their joint law
(Lemma~\ref{lem:gauss_tail}): $\E\max_{i\le m}Z_i \le \sigma\sqrt{2\log m}$.
\end{lemma}

\begin{proof}
\uses{lem:gauss_tail}
Integrability: $\abs{Z_i}\le e^{Z_i} + e^{-Z_i}$ is integrable ($\beta = \pm1$), and
$\abs{\max_iZ_i} \le \sum_i\abs{Z_i}$. Let $\beta > 0$. By Jensen's inequality for the convex
function $\exp$ (Mathlib \texttt{ConvexOn.map\_integral\_le} with \texttt{convexOn\_exp}; the
integrand $e^{\beta\max_iZ_i} \le \sum_ie^{\beta Z_i}$ is integrable by assumption),
\[
  \exp\Big(\beta\,\E\max_iZ_i\Big) \le \E\exp\Big(\beta\max_iZ_i\Big)
  = \E\max_i e^{\beta Z_i} \le \sum_{i\le m}\E e^{\beta Z_i}
  \le m\,e^{\beta^2\sigma^2/2},
\]
where $\max_ie^{\beta Z_i} = e^{\beta\max_iZ_i}$ because $\exp$ is increasing. Taking
logarithms, $\E\max_iZ_i \le \frac{\log m}{\beta} + \frac{\beta\sigma^2}{2}$ for every
$\beta>0$. If $m\ge2$ and $\sigma>0$, choose $\beta := \sqrt{2\log m}/\sigma$: both terms
equal $\sigma\sqrt{2\log m}/2$. If $m = 1$, let $\beta\downarrow0$: $\E Z_1\le0 = \sigma\sqrt{2\log1}$.
If $\sigma = 0$, let $\beta\to\infty$: $\E\max_iZ_i\le0$. The Gaussian case: by
Lemma~\ref{lem:gauss_tail} (with $\mu = 0$), $Z_i\sim\Normal(0,\sigma_i^2)$ with $\sigma_i>0$
has a sub-Gaussian MGF with constant $\sigma_i^2\le\sigma^2$, hence with constant $\sigma^2$
(the bound $e^{\beta^2\sigma_i^2/2}\le e^{\beta^2\sigma^2/2}$ is monotone in the constant);
if $\sigma_i = 0$ then $Z_i = 0$ a.s.\ and the constant-$\sigma^2$ bound is trivial.
\end{proof}

\begin{lemma}[Comparison by covariance domination]\label{lem:gauss_comparison}
\lean{ProbabilityTheory.gaussianWidth_multivariateGaussian_le, ProbabilityTheory.gaussianWidth_le_gaussianWidth_conv}\leanok
\uses{def:pg_gaussian_vector, lem:sup_countable_dense}
Let $\Xs \subset \R^d$ be nonempty and compact, let $S_X, S_Y \in \PSD$ with
$S_X \preceq S_Y$ (i.e.\ $S_Y - S_X \in \PSD$), and let $G_X \sim \Normal(0,S_X)$,
$G_Y \sim \Normal(0,S_Y)$ (on arbitrary probability spaces). Then
$\sup_{x\in\Xs}\ip{x}{G_X}$ and $\sup_{x\in\Xs}\ip{x}{G_Y}$ are integrable and
\[
  \E\sup_{x\in\Xs}\ip{x}{G_X} \le \E\sup_{x\in\Xs}\ip{x}{G_Y}.
\]
In particular (finite index set $\Xs = \{e_1,\dots,e_m\}$, the canonical basis of $\R^m$):
if $Z \sim \Normal(0,\Sigma_Z)$ and $W \sim \Normal(0,\Sigma_W)$ in $\R^m$ with
$\Sigma_Z \preceq \Sigma_W$, then $\E\max_{i\le m}Z_i \le \E\max_{i\le m}W_i$.
\end{lemma}

\begin{proof}
\leanok
\uses{lem:sup_countable_dense, lem:gauss_norm_moments, lem:pg_indep_gaussian_pair, lem:gauss_law_of_cov}
Write $m(v) := \sup_{x\in\Xs}\ip{x}{v}$, which is continuous with $\abs{m(v)} \le R\norm{v}$,
$R := \max_{x\in\Xs}\norm{x}$ (Lemma~\ref{lem:sup_countable_dense}); integrability of $m(G_X)$
and $m(G_Y)$ follows from $\E\norm{G} < \infty$ for Gaussian vectors
(Lemma~\ref{lem:gauss_norm_moments}(iii)). Both expectations depend only on the laws, so we may
work on the product space $(\R^d\times\R^d, \Normal(0,S_X)\otimes\Normal(0,S_Y-S_X))$ with
coordinate projections $G_X$ and $G_Z$, which are independent with
$G_Z \sim \Normal(0, S_Y - S_X)$ ($S_Y - S_X \in \PSD$ by assumption). By
Lemma~\ref{lem:pg_indep_gaussian_pair}, $G_X + G_Z \sim \Normal(0, S_X + (S_Y - S_X)) =
\Normal(0,S_Y)$, so $\E m(G_Y) = \E m(G_X + G_Z)$ (Lemma~\ref{lem:gauss_law_of_cov} is not
even needed: this is equality of laws). By Fubini on the product measure
(Mathlib \texttt{MeasureTheory.integral\_prod}; the integrand $(u,z)\mapsto m(u+z)$ is
continuous and dominated by $R(\norm{u}+\norm{z})$, which is integrable on the product),
\[
  \E m(G_X + G_Z) = \int\Big(\int m(u+z)\,\Normal(0,S_Y-S_X)(dz)\Big)\Normal(0,S_X)(du).
\]
Fix $u$. For every $x \in \Xs$ and every $z$, $\ip{x}{u} + \ip{x}{z} \le m(u+z)$; integrating in
$z$ and using $\int\ip{x}{z}\,\Normal(0,S_Y-S_X)(dz) = \ip{x}{0} = 0$ (centered), monotonicity
of the integral gives $\ip{x}{u} \le \int m(u+z)\,\Normal(0,S_Y-S_X)(dz)$. Taking the supremum
over $x\in\Xs$: $m(u) \le \int m(u+z)\,\Normal(0,S_Y-S_X)(dz)$. Integrating in $u$ yields
$\E m(G_X) \le \E m(G_X+G_Z) = \E m(G_Y)$. The finite-index statement is the case
$\Xs = \{e_1,\dots,e_m\}$, since $\ip{e_i}{G} = G_i$.
\end{proof}

\begin{remark}[Relation with Sudakov--Fernique]\label{rem:pg_comparison_vs_sf}
The Sudakov--Fernique inequality (Chapter~\ref{chap:pre_sudakov}) only requires
$\E(Z_i - Z_j)^2 \le \E(W_i - W_j)^2$ for all pairs, which is weaker than
$\Sigma_Z \preceq \Sigma_W$. The Loewner-order version above is all that the monotonicity of
the Gaussian width (Lemma~\ref{lem:width_mono}) needs, and its proof is elementary; this is
one of the deviations from the paper listed in the introduction.
\end{remark}
